\chapter{svnfsfs Reference---Subversion FSFS Filesystem
Tool}\label{svn.ref.svnfsfs}

\texttt{svnfsfs} is a low-level tool for examining and tuning the FSFS
data store that backs a Subversion repository. It was introduced in
Subversion 1.9 to hold FSFS-specific functionality that does not belong
in the general-purpose \texttt{svnadmin} tool. Most administrators will
never need it; it is intended for diagnostics, for advanced tuning, and
for recovering a damaged repository under the guidance of the Subversion
development community.

Like \texttt{svnadmin}, \texttt{svnfsfs} works via direct repository
access and so must be run on the machine that holds the repository,
referring to it by a local path rather than a URL. Its
\texttt{dump-index} and \texttt{load-index} subcommands operate on the
FSFS logical addressing index and therefore require a format 7 (Subversion
1.9 or later) repository.

\section{svnfsfs dump-index}\label{svn.ref.svnfsfs.c.dump-index}

\emph{Dump the index contents for a revision or pack file.}

\textbf{Synopsis}

\texttt{svnfsfs\ dump-index\ REPOS\_PATH\ -r\ REV}

\subsection{Description}

Print, to stdout, the contents of the logical-addressing index for the
revision or pack file that contains revision \texttt{REV}. This is
available only for format 7 (Subversion 1.9 and later) repositories. The
output begins with a header line and then lists one entry per line,
ordered by location within the revision or pack file, giving each
item\textquotesingle s byte offset, length, type (for example,
\texttt{frep} for a file representation, \texttt{drep} for a directory
representation, \texttt{node} for a node-revision record, and
\texttt{chgs} for a changed-paths list), revision, item number, and
checksum. The output can be edited and fed back in with
\texttt{svnfsfs\ load-index} to repair a corrupted index.

\subsection{Options}

\begin{verbatim}
--memory-cache-size (-M) ARG
--revision (-r) ARG
\end{verbatim}

\subsection{Examples}

\begin{verbatim}
$ svnfsfs dump-index /var/svn/repos -r 1
       Start       Length Type   Revision     Item Checksum
           0           2f frep          1        3 7ee0088a
          2f           1f frep          1        4 3ba36966
          4e           a7 node          1        5 441a8d27
          f5           3b fprop         1        6 7033638d
         130           e3 node          1        7 1efba051
         213           52 drep          1        8 f54298a5
         …
$
\end{verbatim}

\section{svnfsfs help (h, ?)}\label{svn.ref.svnfsfs.c.help}

\emph{Help!}

\textbf{Synopsis}

\texttt{svnfsfs\ help\ {[}SUBCOMMAND...{]}}

\subsection{Description}

This subcommand is useful when you\textquotesingle re trapped on a
desert island with neither a Net connection nor a copy of this book.

\subsection{Options}

None

\section{svnfsfs load-index}\label{svn.ref.svnfsfs.c.load-index}

\emph{Rewrite the index from a dumped-index stream.}

\textbf{Synopsis}

\texttt{svnfsfs\ load-index\ REPOS\_PATH}

\subsection{Description}

Read index contents from stdin and rewrite the FSFS logical-addressing
index accordingly. The input uses the same format produced by
\texttt{svnfsfs\ dump-index}, except that the header and per-entry
checksums are optional and ignored. The data must cover the entire
revision or pack file; the revision number is taken automatically from
the input, and the entries need not be in any particular order. This
subcommand is used to repair a damaged index, typically with input
derived from a \texttt{dump-index} of the affected revision.

\subsection{Options}

\begin{verbatim}
--memory-cache-size (-M) ARG
\end{verbatim}

\subsection{Examples}

Rebuild the index of a revision from a (possibly hand-corrected) dump:

\begin{verbatim}
$ svnfsfs dump-index /var/svn/repos -r 42 > r42.index
$ # ...edit r42.index as needed...
$ svnfsfs load-index /var/svn/repos < r42.index
$
\end{verbatim}

\section{svnfsfs stats}\label{svn.ref.svnfsfs.c.stats}

\emph{Print FSFS storage statistics.}

\textbf{Synopsis}

\texttt{svnfsfs\ stats\ REPOS\_PATH}

\subsection{Description}

Write detailed statistics about the FSFS data store to stdout: how many
bytes are consumed by revisions, changed-path lists, node-revision
records, and representations; how effective representation sharing and
deltification have been; and per-item-type breakdowns. These numbers are
useful for understanding where a repository\textquotesingle s on-disk
size comes from and for gauging the benefit of packing or repacking.

\subsection{Options}

\begin{verbatim}
--memory-cache-size (-M) ARG
\end{verbatim}

\subsection{Examples}

\begin{verbatim}
$ svnfsfs stats /var/svn/repos

Global statistics:
               1,868 bytes in            3 revisions
                 203 bytes in            4 changes
               1,225 bytes in            8 node revision records
                 316 bytes in            9 representations
                 292 bytes expanded representation size
                 326 bytes with rep-sharing off
…
$
\end{verbatim}
